\cotefile{29}
\scanpage{26}%
% archive sheet number: 25
have in addition a \uline{fiber functor}
\[
F:\operatorname{Rev}^{\underline{R}}(S)\qquad(\mathrm{Ens\ finite}),
\]
i.e. of an exact functor\inserted{(conditions G4,G5 of V 4)}. In the case of example 1, every geometric point of $Z$ provides such a fiber functor, and this functor is conservative (condition G6) if one assumes that every open and closed subset of $Z$ stable under $G$ is empty or equal to $Z$. Likewise, in the case of example 2, one obtains a \deleted{\unsure{pai|by}} fiber functor by taking a geometric point of $Z=\lim Z_i$. To obtain a fiber functor in the more general case of example 3, one chooses a geometric point $\qquad S$ of $S$ \deleted{such that} which does not belong to\deleted{\unsure{u|n}} ``the ramification set''of $\underline{R}$, i.e. such that for every $S'$ in $\mathrm{Et}_S$ and every object $M$ of $\underline{C}_{S'}$, \unsure{\deleted{'}the|\deleted{e}the} \deleted{object} \inserted{covering with} \deleted{with} operators $\underline{r}(M)$ \deleted{over $S'$}\inserted{(group $G$)} of $S'$ is principal at the points of $S'$above the the image \inserted{$s$} of $\quad$ in $S$. Then \deleted{for every étale covering} \deleted{$T$ with group $G$ \unsure{of the|d} of \unsure{to a|of a|the one} of the scheme with operators $G$} one observes that for every object $T$ of $\operatorname{Rev}^{\underline{R}}(S)$, $\underline{r}(T)$ is étale enr $S$ at $s$, and the functor $T\quad\underline{r}(T)(\quad)$ (geometric fiber at $\quad$) is indeed exact. For it to be conservative, it is necessary and sufficient that the site $\underline{C}$ (for the inverse image topology of that of $\underline{\mathrm{Et}}_S$) be connected; \deleted{i.e. that every covering sieve \unsure{of the|of} this \unsure{site|site}} if one assumes, with the notation of example 3, that for every $M$, defining $(\underline{r}(M),G)$, $\underline{r}(M)/G\qquad S'$ is an isomorphism (which will probably be the case in all interesting cases), then it amounts to the same thing to say that $S$ is connected. Applying the theory of V 4, one thus finds a profinite group, denoted ${}^{\underline{R}}_1(S,\quad)$, \deleted{which classifi} and an equivalence of the category of finite sets on which this group acts continuously, with the category $\operatorname{Rev}^{\underline{R}}(S)$.

3. \uline{Functorialities for} $S$ \inserted{\uline{and} $\underline{R}$} \uline{varying}

\unsure{Consider|Consider} a morphism $f:T\qquad S$ of schemes, whence a functor
